\section{Resource Mapping}\label{sec:methodology}

LLM evaluation covers more than checking factual answers: it also examines
whether a response follows an intended request, uses tools, develops creative
content, or solves a constrained problem. We compare forty-one resources
across these task families to give a broader view of how evaluation designs
handle reference evidence, permission, required status marking, and style.
The comparison includes datasets, benchmarks, and scoring protocols, which
we collectively call \emph{evaluation resources}.

The resulting \emph{resource mapping} connects the framework to a range of
evaluation settings. Its categories provide a way to compare task designs
beyond any one dataset; findings about their prevalence remain limited to
the selected inventory. The mapping does not test the labeling rule's
reliability in practice.

\subsection{Selection and Scope}

The unit of analysis is an \emph{evaluation resource}: a benchmark, dataset,
or protocol that specifies how responses are evaluated. The six task categories listed in
Tables~\ref{tab:full-mapping-factual-a}--\ref{tab:full-mapping-creative} cover factuality evaluation, intent evaluation,
mixed-task and tool-use evaluation, creative writing and reasoning, divergent
thinking and linguistic creativity, and constrained problem solving.
The inventory brings together hallucination-related and creativity-related
resources identified for this paper, extending the comparison beyond
factual question answering. It supports a general discussion of evaluation
designs, without establishing exhaustive coverage of the field. We call
this selection a \emph{purposive sample}.

Tables~\ref{tab:full-mapping-factual-a}--\ref{tab:full-mapping-creative}
list all 41 resources, their references, and the evaluation criteria in the
main text. The available record does not explain how the resources were
initially found or selected. That procedure must be documented for others
to reproduce the selection; the final list alone is insufficient.

The source check was uneven: some entries were examined through relevant
methodological passages, others through primary summaries, and one source
was only partly accessible. These limits affect how confidently we can
interpret each protocol. In particular, failing to identify a criterion
in the reviewed material does not establish that the resource lacks it.
Recent designs discussed in Section~\ref{sec:related} remain outside the
41-resource inventory until the same comparison has been completed.
% AUTHOR ACTION: replace the preceding limitation with the actual documented
% discovery procedure when available. Do not invent dates, queries, or screening.

\subsection{Coding Protocol}

For each resource, we assign categories to describe its evaluation design.
This process, called \emph{coding}, records seven aspects in three groups:
\begin{itemize}
\item Contract information: the task oracle type, permission scope, and required
status marking.
\item Task information: the type of task and any separate criterion for
quality or task success.
\item Style evaluation: how the requested style is treated, and whether
claims or labels are compared across different wordings that preserve claim
content and observed status marking.
\end{itemize}
We group the resources into six task categories, using the same names
throughout the paper: \emph{factuality evaluation} (IDs 1--22),
\emph{intent evaluation} (23), \emph{mixed-task and tool-use evaluation}
(24--27), \emph{creative writing and reasoning} (28--32, 35--36),
\emph{divergent thinking and linguistic creativity} (33--34), and
\emph{constrained problem solving} (37--41).

Factuality evaluation checks answers or summaries against evidence. Intent
evaluation checks interpretation of a request. Mixed-task and tool-use
evaluation combines subtasks or interactions. Creative writing and reasoning
assesses stories, written explanations, or alternative reasoning paths.
Divergent thinking and linguistic creativity assesses the variety and
originality of ideas or language. Constrained problem solving checks whether
a proposed solution meets stated requirements. These categories summarize
the task and evaluation criteria recorded for each resource; they do not
rank resources or determine whether an individual claim is permitted.
They are descriptive groupings, not an independently validated taxonomy.
Figure~\ref{fig:profiles} uses these same six categories.

The descriptions reflect our reading of each resource, whose authors may
use different terms. A score serves as reference evidence only if it is used
to support or contradict the evaluated content; measuring overall task
success alone is not enough.

\subsection{Resource Tables}\label{subsec:resource-tables}

Read each row from left to right: the resource and its citation; the task
and the evidence or checks used; other criteria such as quality, novelty,
or feasibility; and our interpretation of what content the task allows.
The category headings use the same names as Figure~\ref{fig:profiles}.

``No extra quality measure identified'' means that this review did not
identify a separate quality criterion beyond the evidence or correctness
check in the preceding column. ``Only supported content (inferred)'' means
that we interpret the task as requiring support for its claims.
``Permission not established by this review'' means that we have not
established which unsupported additions the task permits. It does not
mean that all additions are forbidden or that the resource has no instructions.
The permission descriptions for creative tasks likewise summarize our
interpretation of the task or fictional setting.

Status-marking requirements have not yet been verified across these
resources, so the tables do not assign them a positive or negative result.
Style is discussed separately below: a missing controlled comparison of
style does not imply that a resource ignores writing quality.

\begingroup
\small
\setlength{\tabcolsep}{3pt}
\renewcommand{\arraystretch}{1.0}
\begin{longtable}{@{}P{.20\linewidth}P{.32\linewidth}P{.24\linewidth}P{.19\linewidth}@{}}
\caption{\textbf{Factuality evaluation (resources 1--13).} The final column interprets task instructions; it is not a verified claim-level permission rule.}
\label{tab:full-mapping-factual-a}\\
\toprule
\textbf{Resource} & \textbf{What is checked, and against what?} & \textbf{Other evaluation criteria} & \textbf{Allowed content: our reading}\\\midrule
\endfirsthead
\multicolumn{4}{l}{\textbf{Table \thetable{} (continued)}}\\
\toprule
\textbf{Resource} & \textbf{What is checked, and against what?} & \textbf{Other evaluation criteria} & \textbf{Allowed content: our reading}\\\midrule
\endhead
\midrule\multicolumn{4}{r}{\emph{Continued on next page}}\\
\endfoot
\bottomrule
\endlastfoot
\multicolumn{4}{@{}p{.96\linewidth}@{}}{\textbf{Task family: factuality evaluation}}\\*
1.~FEVER \citep{thorne_fever_2018} & Fact checking\newline Wikipedia evidence sentences. & No extra quality measure identified & Only supported content (inferred)\\\addlinespace[3pt]
2.~TruthfulQA \citep{lin_truthfulqa_2022} & Factual QA\newline Sourced reference answers; human adjudication for generations. & Informativeness & Only supported content (inferred)\\\addlinespace[3pt]
3.~PopQA \citep{mallen_when_2023} & Long-tail QA\newline Wikidata. & No extra quality measure identified & Only supported content (inferred)\\\addlinespace[3pt]
4.~SimpleQA \citep{wei_measuring_2024} & Factual QA\newline Verified reference answers. & Answer correctness and not-attempted outcome & Only supported content (inferred)\\\addlinespace[3pt]
5.~TriviaQA \citep{joshi_triviaqa_2017} & Reading QA\newline Gold answers plus Web/Wikipedia documents. & No extra quality measure identified & Only supported content (inferred)\\\addlinespace[3pt]
6.~DefAn \citep{rahman_defan_2024} & Factual QA\newline Definitive reference answers. & Prompt alignment and cross-run consistency & Only supported content (inferred)\\\addlinespace[3pt]
7.~HaluEval \citep{li_halueval_2023} & Multi-task hallucination recognition\newline Task references or source text; Web-assisted adjudication. & Specificity and relevance folded into some labels & Only supported content (inferred)\\\addlinespace[3pt]
8.~FactBench \citep{bayat_factbench_2024} & Open-ended factuality evaluation\newline Web-retrieved material. & No extra quality measure identified & Only supported content (inferred)\\\addlinespace[3pt]
9.~HalluLens \citep{bang_hallulens_2025} & Extrinsic/intrinsic hallucination benchmark\newline Wikipedia, nonexistent entities, or task context. & Correctness and refusal/acceptance & Only supported content (inferred)\\\addlinespace[3pt]
10.~FRANK \citep{pagnoni_understanding_2021} & Summarization\newline Source documents. & No extra quality measure identified & Only supported content (inferred)\\\addlinespace[3pt]
11.~QAGS \citep{wang_asking_2020} & Summarization\newline Source documents. & No extra quality measure identified & Only supported content (inferred)\\\addlinespace[3pt]
12.~SummEval \citep{fabbri_summeval_2021} & Summarization\newline Source documents. & Automatic metrics and human summarization-quality judgments & Only supported content (inferred)\\\addlinespace[3pt]
13.~XSum \citep{narayan_xsum_2018} & Summarization\newline BBC source article. & ROUGE, human preference, and information preservation & Only supported content (inferred)\\
\end{longtable}
\endgroup

\begingroup
\small
\setlength{\tabcolsep}{3pt}
\renewcommand{\arraystretch}{1.0}
\begin{longtable}{@{}P{.20\linewidth}P{.32\linewidth}P{.24\linewidth}P{.19\linewidth}@{}}
\caption{\textbf{Factuality evaluation, intent evaluation, and mixed-task and tool-use evaluation (resources 14--27).} The final column interprets task instructions; it is not a verified claim-level permission rule.}
\label{tab:full-mapping-factual-b}\\
\toprule
\textbf{Resource} & \textbf{What is checked, and against what?} & \textbf{Other evaluation criteria} & \textbf{Allowed content: our reading}\\\midrule
\endfirsthead
\multicolumn{4}{l}{\textbf{Table \thetable{} (continued)}}\\
\toprule
\textbf{Resource} & \textbf{What is checked, and against what?} & \textbf{Other evaluation criteria} & \textbf{Allowed content: our reading}\\\midrule
\endhead
\midrule\multicolumn{4}{r}{\emph{Continued on next page}}\\
\endfoot
\bottomrule
\endlastfoot
\multicolumn{4}{@{}p{.96\linewidth}@{}}{\textbf{Task family: factuality evaluation (continued)}}\\*
14.~CNN/\allowbreak DailyMail \citep{hermann_teaching_2015,wang_asking_2020} & News corpus used in summarization evaluation\newline News source article. & No extra quality measure identified & Only supported content (inferred)\\\addlinespace[3pt]
15.~DelucionQA \citep{sadat_delucionqa_2023} & Domain QA\newline Retrieved chunks from one Jeep Gladiator manual. & Answerability, coverage, and abstention & Only supported content (inferred)\\\addlinespace[3pt]
16.~HalluQA \citep{cheng_evaluating_2023} & Chinese adversarial QA\newline Answer sets backed by sources or explanations. & Answer correctness and refusal handling & Only supported content (inferred)\\\addlinespace[3pt]
17.~CR-LLM \citep{li_cr-llm_2024} & Concept reasoning\newline Reference concepts plus GPT-4 contextual adjudication. & Concept granularity & Only supported content (inferred)\\\addlinespace[3pt]
18.~HotpotQA \citep{yang_hotpotqa_2018} & Multi-hop QA\newline Wikipedia evidence passages. & No extra quality measure identified & Only supported content (inferred)\\\addlinespace[3pt]
19.~HalluRAG \citep{ridder_hallurag_2025} & RAG hallucination detection\newline Provided passages/answers and derived sentence labels. & No extra quality measure identified & Only supported content (inferred)\\\addlinespace[3pt]
20.~RealTime QA \citep{kasai_realtime_2024} & Real-time QA\newline Time-indexed reference answers. & No extra quality measure identified & Only supported content (inferred)\\\addlinespace[3pt]
21.~HypoTermQA \citep{uluoglakci_hypotermqa_2024} & Hypothetical-term evaluation\newline Term-validity labels and Wikipedia definitions. & No extra quality measure identified & Only supported content (inferred)\\\addlinespace[3pt]
22.~AA-Omniscience \citep{jackson_aaomniscience_2025,artificialanalysis_aaomniscience_public_2025} & Cross-domain factual QA\newline Source-derived reference answers, evaluated by an LLM grader. & Accuracy, incorrect-answer penalty, and abstention-sensitive reliability & Only supported content (inferred)\\\addlinespace[3pt]
\midrule
\multicolumn{4}{@{}p{.96\linewidth}@{}}{\textbf{Task family: intent evaluation}}\\*
23.~FaithQA \citep{hao_faithqa_2025} & Intent verification\newline RAG reports; Wikipedia for factual QA; none for creative writing. & Query-constraint satisfaction & Permission not established by this review\\\addlinespace[3pt]
\midrule
\multicolumn{4}{@{}p{.96\linewidth}@{}}{\textbf{Task family: mixed-task and tool-use evaluation}}\\*
24.~BIG-Bench \citep{srivastava_beyond_2023} & Multi-task reasoning\newline Subtask targets and adjudication. & No extra quality measure identified & Permission not established by this review\\\addlinespace[3pt]
25.~BIG-Bench Hard \citep{suzgun_challenging_2023} & Reasoning suite\newline Task-specific reference answers. & No extra quality measure identified & Permission not established by this review\\\addlinespace[3pt]
26.~ToolQA \citep{zhuang_toolqa_2023} & Tool-use QA\newline Fixed corpora and generated gold answers. & No extra quality measure identified & Permission not established by this review\\\addlinespace[3pt]
27.~HIC-Bench \citep{yang_hicbench_2025} & Hallucination typing\newline CDKB/Wikipedia principles and human-verified judgments. & Originality, feasibility, value, and composite metrics & Permission not established by this review\\
\end{longtable}
\endgroup

\begingroup
\small
\setlength{\tabcolsep}{3pt}
\renewcommand{\arraystretch}{1.0}
\begin{longtable}{@{}P{.20\linewidth}P{.32\linewidth}P{.24\linewidth}P{.19\linewidth}@{}}
\caption{\textbf{Creative writing and reasoning, divergent thinking and linguistic creativity, and constrained problem solving (resources 28--41).} The final column interprets task instructions; it is not a verified claim-level permission rule.}
\label{tab:full-mapping-creative}\\
\toprule
\textbf{Resource} & \textbf{What is checked, and against what?} & \textbf{Other evaluation criteria} & \textbf{Allowed content: our reading}\\\midrule
\endfirsthead
\multicolumn{4}{l}{\textbf{Table \thetable{} (continued)}}\\
\toprule
\textbf{Resource} & \textbf{What is checked, and against what?} & \textbf{Other evaluation criteria} & \textbf{Allowed content: our reading}\\\midrule
\endhead
\midrule\multicolumn{4}{r}{\emph{Continued on next page}}\\
\endfoot
\bottomrule
\endlastfoot
\multicolumn{4}{@{}p{.96\linewidth}@{}}{\textbf{Task family: creative writing and reasoning}}\\*
28.~CreativityPrism \citep{hou_creativityprism_2025} & Multi-task creativity\newline Tests or correctness judges where applicable. & Task-specific rubrics & Invention within task instructions\\\addlinespace[3pt]
29.~LitBench \citep{fein_litbench_2025} & Story evaluation\newline No external factual reference identified (human preferences assess quality). & Human preferences & Fictional content within the story\\\addlinespace[3pt]
30.~Short Story \citep{ismayilzada_evaluating_2025} & Story generation\newline No external factual reference identified (cue words constrain the task, not claim truth). & Novelty, diversity, complexity, and judged creativity & Fictional content within the story\\\addlinespace[3pt]
31.~WritingBench \citep{wu_writingbench_2025} & Generative writing\newline Provided materials when present. & Writing quality and requested style & Invention within task instructions\\\addlinespace[3pt]
32.~Art-or-Artifice / TTCW \citep{chakrabarty_art_2024} & Creative writing\newline No external factual reference identified (expert tests assess creativity). & Expert rubric & Fictional content within the story\\\addlinespace[3pt]
35.~Automated Creativity \citep{li_automated_2025} & Creative writing\newline No external factual reference identified (reference texts anchor quality, not claim truth). & Reference-text quality anchor & Invention within task instructions\\\addlinespace[3pt]
36.~Hallucinating LLM Could Be Creative \citep{sun_hallucinating_2024} & Reasoning-path evaluation\newline Task correctness; consistency/clustering are diagnostics. & Correctness, consistency, diversity & Permission only partly characterized\\\addlinespace[3pt]
\midrule
\multicolumn{4}{@{}p{.96\linewidth}@{}}{\textbf{Task family: divergent thinking and linguistic creativity}}\\*
33.~Automated DT scoring \citep{organisciak_beyond_2023} & Divergent thinking\newline No external factual reference identified. & Novelty and variety of ideas & New ideas within the task\\\addlinespace[3pt]
34.~Verbal Creativity \citep{dinu_testing_2025} & Linguistic creativity\newline No external factual reference identified. & Originality, flexibility, elaboration & Creative language within the task\\\addlinespace[3pt]
\midrule
\multicolumn{4}{@{}p{.96\linewidth}@{}}{\textbf{Task family: constrained problem solving}}\\*
37.~MacGyver \citep{tian_macgyver_2025} & Problem solving\newline Human judgment of feasibility and efficiency. & Feasibility and efficiency & Permission not established by this review\\\addlinespace[3pt]
38.~EscapeBench \citep{qian_escapebench_2025} & Agent puzzles\newline Game state and progress. & Feasibility or correctness & Permission not established by this review\\\addlinespace[3pt]
39.~CreativEval \citep{delorenzo_creativeval_2024} & Code generation\newline Functionality under Verilog testbenches. & Fluency, flexibility, originality, elaboration & Permission not established by this review\\\addlinespace[3pt]
40.~NeoCoder \citep{lu_benchmarking_2025} & Code generation\newline Test success and constraint adherence. & Correctness, constraint adherence, divergence & Permission not established by this review\\\addlinespace[3pt]
41.~CreativeMath \citep{ye_assessing_2025} & Mathematical solution generation\newline Reference solutions and unanimous LLM correctness judgment. & Correctness and solution novelty & Permission not established by this review\\
\end{longtable}
\endgroup


\subsection{Resource Analysis Criteria}\label{sec:mapping}

\paragraph{Reference evidence and reported checks}
The inventory records source documents, reference answers, retrieved
material, and task-validity checks. These are different bases for evaluation.
For example, FEVER uses evidence sentences (ID~1), QAGS compares answers
derived from a summary and its source (11), and SimpleQA uses reference
answers with a grading procedure (4)
\citep{thorne_fever_2018,wang_asking_2020,wei_measuring_2024}.
Example~A in Section~\ref{sec:examples} illustrates why a label must retain the
identity of the reference used.

For constrained problem solving, the inspected protocols evaluate solution
validity: feasibility under stated constraints in MacGyver (37), progress
in the game environment in EscapeBench (38), functionality under Verilog
testbenches in CreativEval (39), test success and constraint adherence in
NeoCoder (40), and LLM-judged mathematical correctness in CreativeMath (41)
\citep{tian_macgyver_2025,qian_escapebench_2025,delorenzo_creativeval_2024,
lu_benchmarking_2025,ye_assessing_2025}. These checks do not automatically
verify every claim in an accompanying explanation. The tables therefore
report their stated role, rather than count them as verified implementations
of a single claim-level oracle.

For LitBench, Short Story, Art-or-Artifice, Automated DT scoring, Verbal
Creativity, and Automated Creativity (IDs 29, 30, 32--35), no external
factual reference was identified in the inspected evaluation criteria.
Their prompts, story requirements, quality references, or creativity
rubrics still constrain evaluation. This observation does not establish
the absence of all evidence checks or all task constraints.

\paragraph{Permission scope}
We interpret factuality evaluation tasks as allowing supported content
only (22 resources, IDs 1--22). For eight creative resources (28--35),
permission is inferred from the task or fictional setting. For resource 36,
the permission conditions are only partly characterized. For the remaining
ten resources (23--27 and 37--41), this review has not established the scope
of permitted additions. These are interpretations of the selected tasks,
not verified claim-level permission rules. Section~\ref{sec:agenda}
describes the protocol audit needed to check them against the full
contract definition.

\paragraph{Required status marking}
The status-marking field is \emph{unverified} throughout this inventory:
the present source check has not established, for each row, both a
requirement on how a claim's status must be presented and the procedure
that verifies it. This is the status of our audit, not a finding that the
resources lack such requirements. We therefore do not report counts that
contrast absent and embedded marking checks.

Several resources explicitly assess abstention or refusal. SimpleQA
includes a not-attempted outcome, HalluLens evaluates refusal, and
DelucionQA records answerability and whether a response answers
\citep{wei_measuring_2024,bang_hallulens_2025,sadat_delucionqa_2023}.
AA-Omniscience also accounts for abstention in reliability evaluation
\citep{jackson_aaomniscience_2025}. These measures concern whether an answer
is offered. They do not by themselves establish how a permitted hypothesis
or fictional claim must be marked under a truth contract.

\subsection{Joint Evaluation and Style Variation}

WritingBench explicitly scores writing requirements, including style
\citep{wu_writingbench_2025}. Our inventory records this separately from a
comparison that holds claim content and observed status marking fixed.
No such controlled comparison is documented in the current coding. That
statement describes the inventory; a comprehensive claim of absence would
require inspecting every relevant protocol and subtask.

Scoring requested style asks whether a response uses the style requested
by the task. A comparison with fixed claims asks a different question:
does changing only the style change the evaluator's judgment? Our review
has not documented such a comparison for the selected resources. For action
sequences or executable outputs, a test of prose style may not be applicable.

Figure~\ref{fig:profiles} groups the selected resources by task family and
places the joint evaluation sought by this paper at the center: assessing
creativity and divergent thinking while distinguishing licensed divergence
from hallucination under an explicit truth contract. The groups summarize
the resource tables; the central area identifies a research target, rather
than a measured empty region in an exhaustive survey.

Some selected resources already combine creativity-related criteria with
correctness checks or hallucination typing, as recorded for HIC-Bench and
constrained problem solving in Table~\ref{tab:full-mapping-creative}.
The question is whether such evaluations distinguish permitted invention
from hallucination using reference evidence, permission scope, and required
status marking. The current coding does not establish that no existing
protocol does so; that absence claim requires source-level verification.
Style variation remains a separate test of whether an evaluator's judgment
changes when claim content and observed status marking stay fixed.

\begin{figure}[!ht]
\centering
\resizebox{\textwidth}{!}{\input{figures/resource_clusters.tex}}
\caption{\textbf{Resource families and the proposed evaluation target.}
Colored clusters group selected resources by task type; examples and resource
IDs identify each group. The six category names match the tables in
Section~\ref{subsec:resource-tables}.
The central area states the evaluation sought: creativity and divergent
thinking assessed together with hallucination detection under an explicit
truth contract. Related combinations already occur in the selection; the
figure does not establish that this target is absent from all prior work.
Positions, distances, and cluster sizes carry no quantitative meaning.
Resource-level comparisons appear in
Tables~\ref{tab:full-mapping-factual-a}--\ref{tab:full-mapping-creative}.}
\label{fig:profiles}
\end{figure}
\FloatBarrier

